% page 201 of the facsimile (image p-200 of the scan)
% block 167.6 x 246.3 mm ; left margin 34.4 mm
\pgc{12.700mm}{0.000mm}{
\l{}
\l{}
\l{}
\l{}
\l{\cel{58}193}
\l{}
\l{\sou{gibelino}. (\sou{historio di Italia}). Dum la Mezepoko, partisano di}
\l{\cel{3}la imperiestro Germana, opoze a la guelfi, partisani di la}
\l{\cel{3}papo. - DEFIRS.}
\l{}
\l{\sou{gibono}.\cel{2}Granda simio di India. - DEF.}
\l{}
\l{\sou{gicheto}. Mikra aperturo tra qua onu komunikas kun la oficieri}
\l{\cel{3}en la kontori postala, kun la kasieri en la firmi banka e}
\l{\cel{3}komercala. - F.}
\l{}
\l{\sou{giganto}. Persono di qua la staturo esas extraordinare longa.}
\l{\cel{3}- DEFIRS.}
\l{}
\l{\sou{gildo}. Ensemblo de sama-speca mestieristi. - DER.}
\l{}
\l{\sou{gilochar}. (\sou{trans}.)\cel{2}(\sou{tekn.}) Ornar per interplekturo de}
\l{\cel{3}traiti grabita. - DEFR.}
\l{}
\l{\sou{gilotino}. Instrumento por tranchar la kapo di la personi kondam-}
\l{\cel{3}nita ad ocideso legala - ek lamo tre pezoza qua, glitante}
\l{\cel{3}inter du fosti, falas bruske sur la kolo di la ocidoto sternita}
\l{\cel{3}sub olu. - DEFIRS.}
\l{}
\l{\sou{gimbardo}. (\sou{muziko)}. Muzik-instrumento, ek mi-cirklo stala o la-}
\l{\cel{3}tuna qua finas per du branchi paralela, inter qui esas langeto}
\l{\cel{3}stala quan onu agitas per la fingro dum ke onu tenas la instru-}
\l{\cel{3}mento inter la maxilo e la mandibulo.}
\l{}
\l{\sou{gimnastikar}. (\sou{netrans.}) I. Exekutar la exerco-movi apta plufort-}
\l{\cel{3}igar e plu-flexebligar la korpo. - II. Exekutar la exerco-}
\l{\cel{3}taski apta plu-fortigar e plu-flexebligar la fakultati pensala.}
\l{\cel{3}- DEFIRS.}
\l{}
\l{gimnazio. Nomo donata, en Germania, a la kolegii ube onu edukesas}
\l{\cel{3}klasikale. - D.}
\l{}
\l{\sou{gimnospermo}. Di qua la semini aspektas quaze \textquotedbl{}nuda\textquotedbl{}. - DEFIS.}
\l{}
\l{\sou{gimnoto}. (\sou{zool.}) Fisho ek la regioni tropikala di Amerika, sen}
\l{\cel{3}floso dorsala, e qua havas, an omna flanki di la kaudo, lami}
\l{\cel{3}membrana qui emisas elektro. - L.\cel{1}\sou{gymnotua electricus}.- DEFIS.}
\l{}
\l{\sou{ginandra}. (\sou{bot.}) Di qua la stamini insertesas sur la pistilo.-}
\l{\cel{3}DEFIRS.}
\l{}
\l{\sou{gineceo}. I. (\sou{antique}). Parto de la domo qua rezervesis a la}
\l{\cel{3}homini. - II. Loko ube permanis ordinare la homini, la}
\l{\cel{3}mulieri, ube li laboris. - DEFIRS.}
\l{}
\l{\sou{ginesto}. (\sou{bot.}) Arboreto, di qua la floro esas flava, de la}
\l{\cel{3}familio \textquotedbl{}leguminosi\textquotedbl{}. - L.\cel{1}\sou{genista}. - DFIS.}
}
